% ══════════════════════════════════════════════════════════════════════════════
%                         فصل نهم
%                 محافظه‌کاری: از برک تا امروز
% ══════════════════════════════════════════════════════════════════════════════

\chapter{محافظه‌کاری: از برک تا امروز}
\label{ch:conservatism}

\begin{flushright}
\textit{«جامعه قراردادی است میان مردگان، زندگان و آنان که هنوز زاده نشده‌اند.»}

— ادموند برک
\end{flushright}

% ══════════════════════════════════════════════════════════════════════════════
%                         مقدمه فصل
% ══════════════════════════════════════════════════════════════════════════════

\lettrine[lines=3, lraise=0.1, nindent=0.5em]{\textcolor{RightBlue}{م}}{حافظه‌کاری} 
شاید گمراه‌کننده‌ترین نام در سیاست باشد. چه چیزی را «حفظ» می‌کند؟ وضع موجود؟ سنت؟ امتیازات؟ آزادی؟ بسته به زمان و مکان، محافظه‌کاری معانی بسیار متفاوتی داشته است. محافظه‌کار روسی قرن نوزدهم خواستار حفظ تزاریسم بود؛ محافظه‌کار آمریکایی امروز خواستار بازار آزاد است. در این فصل، این تنوع را بررسی می‌کنیم و می‌کوشیم هسته مشترکی در میان شاخه‌های محافظه‌کاری بیابیم.

% ══════════════════════════════════════════════════════════════════════════════
\section{محافظه‌کاری چیست؟}
% ══════════════════════════════════════════════════════════════════════════════

\subsection{تعریف: منفی یا مثبت؟}

\begin{defbox}[تعریف: محافظه‌کاری (رویکرد منفی)]
محافظه‌کاری به مثابه \textbf{مخالفت با تغییر رادیکال}:
\begin{itemize}[nosep]
    \item ترس از انقلاب و آشوب
    \item ترجیح اصلاحات تدریجی بر تغییرات ناگهانی
    \item احتیاط در برابر آزمایش‌های اجتماعی
    \item «اگر خراب نیست، درستش نکن»
\end{itemize}
این تعریف، محافظه‌کاری را صرفاً \textbf{واکنشی} می‌بیند.
\end{defbox}

\begin{defbox}[تعریف: محافظه‌کاری (رویکرد مثبت)]
محافظه‌کاری به مثابه \textbf{دفاع از ارزش‌های خاص}:
\begin{itemize}[nosep]
    \item احترام به \textbf{سنت} و نهادهای تاریخی
    \item اهمیت \textbf{نظم} و ثبات اجتماعی
    \item ارزش \textbf{خانواده}، جامعه محلی و مذهب
    \item شک به \textbf{عقل انتزاعی} و ایده‌های آرمان‌شهری
    \item پذیرش \textbf{نابرابری طبیعی} و سلسله‌مراتب
\end{itemize}
این تعریف، محافظه‌کاری را \textbf{ایدئولوژی مستقل} می‌بیند.
\end{defbox}

\begin{debatebox}[محافظه‌کاری: ایدئولوژی یا temperament؟]
\textbf{محافظه‌کاری به مثابه سرشت (مایکل اوکشات):}
\begin{itemize}[nosep]
    \item محافظه‌کاری یک «گرایش» است، نه مجموعه باورها
    \item ترجیح آشنا بر ناآشنا، موجود بر ممکن
    \item نه راست است نه چپ—رویکردی به زندگی
\end{itemize}

\textbf{محافظه‌کاری به مثابه ایدئولوژی:}
\begin{itemize}[nosep]
    \item محافظه‌کاری مجموعه‌ای از باورهای منسجم دارد
    \item این باورها قابل بیان و دفاع‌اند
    \item احزاب محافظه‌کار برنامه دارند
\end{itemize}

\textbf{شاید هر دو درست باشند:} یک سرشت محافظه‌کارانه هست، اما ایدئولوژی‌های محافظه‌کار مختلفی هم وجود دارند.
\end{debatebox}

\subsection{ویژگی‌های مشترک}

\begin{table}[htbp]
\centering
\caption{ویژگی‌های مشترک اغلب محافظه‌کاران}
\label{tab:conservatism-features}
\begin{tabular}{r|r}
\toprule
\textbf{موضوع} & \textbf{موضع محافظه‌کارانه} \\
\midrule
طبیعت انسان & ناقص، نیازمند کنترل و نهادها \\
تغییر & تدریجی و محتاطانه \\
سنت & مخزن حکمت نسل‌ها \\
عقل & محدود؛ تجربه مهم‌تر از نظریه \\
جامعه & ارگانیک، نه قراردادی \\
دولت & لازم برای نظم، اما محدود \\
مالکیت & حق و مسئولیت \\
مذهب & پایه اخلاق و نظم اجتماعی \\
\bottomrule
\end{tabular}
\end{table}

% ══════════════════════════════════════════════════════════════════════════════
\section{ادموند برک: پدر محافظه‌کاری مدرن}
% ══════════════════════════════════════════════════════════════════════════════

\begin{personbox}[ادموند برک (۱۷۲۹-۱۷۹۷)]

\textbf{ملیت:} ایرلندی-بریتانیایی

\textbf{شغل:} سیاستمدار ویگ، فیلسوف، نویسنده

\textbf{اثر کلیدی:} \book{تأملاتی درباره انقلاب فرانسه} (۱۷۹۰)

\mediumspace

\person{ادموند برک} پارادوکسی جالب است: او ویگ (لیبرال) بود و از انقلاب آمریکا دفاع کرد، اما با انقلاب فرانسه به شدت مخالفت کرد و پدر محافظه‌کاری مدرن شد.

\textbf{استدلال‌های کلیدی برک:}
\begin{enumerate}[nosep]
    \item \textbf{خرد جمعی:} سنت‌ها محصول قرن‌ها آزمون‌وخطا هستند؛ خرد یک فیلسوف نمی‌تواند جای آن‌ها را بگیرد
    \item \textbf{جامعه ارگانیک:} جامعه موجودی زنده است، نه ماشینی که بتوان قطعاتش را عوض کرد
    \item \textbf{قرارداد بین‌نسلی:} جامعه متعلق به مردگان و آیندگان هم هست، نه فقط زندگان
    \item \textbf{پیش‌داوری:} پیش‌داوری‌ها (سنت، عادت) اغلب حکمتی پنهان دارند
    \item \textbf{اصلاح، نه انقلاب:} تغییر باید تدریجی باشد تا نهادها حفظ شوند
\end{enumerate}

\mediumspace

\textbf{نقل‌قول:} \textit{«مردم هرگز نباید یکباره و کامل از گذشته‌شان جدا شوند.»}
\end{personbox}

\begin{quotebox}[ادموند برک، تأملات]
\textit{«ما از عقل فردی خود می‌ترسیم، زیرا گمان می‌کنیم سرمایه عقل در هر انسان اندک است، و بهتر است افراد از انبار عمومی عقل ملت‌ها و قرون بهره ببرند.»}
\end{quotebox}

\begin{quotebox}[ادموند برک، درباره پیش‌داوری]
\textit{«پیش‌داوری در لحظه آماده است؛ پیشاپیش ذهن را در مسیری ثابت از حکمت و فضیلت قرار می‌دهد، و انسان را در لحظه تردید رها نمی‌کند.»}
\end{quotebox}

\subsection{برک و انقلاب: چرا آمریکا آری، فرانسه نه؟}

\begin{comparebox}[مقایسه: برک درباره دو انقلاب]
\begin{table}[H]
\centering
\begin{tabular}{r|c|c}
\toprule
\textbf{موضوع} & \textbf{انقلاب آمریکا} & \textbf{انقلاب فرانسه} \\
\midrule
هدف & حفظ حقوق موجود & ساختن جامعه نو \\
روش & بر اساس سنت انگلیسی & بر اساس ایده‌های انتزاعی \\
رهبران & «اشراف طبیعی» & روشنفکران و عوام \\
نتیجه & نظم جدید پایدار & هرج‌ومرج و ترور \\
موضع برک & حمایت & مخالفت شدید \\
\bottomrule
\end{tabular}
\end{table}
\end{comparebox}

% ══════════════════════════════════════════════════════════════════════════════
\section{محافظه‌کاری قاره‌ای اروپا}
% ══════════════════════════════════════════════════════════════════════════════

در اروپای قاره‌ای، محافظه‌کاری شکل‌های متفاوتی گرفت—اغلب رادیکال‌تر و ضدانقلابی‌تر از نسخه بریتانیایی.

\subsection{ضدانقلابیون فرانسوی}

\begin{personbox}[جوزف دومستر (۱۷۵۳-۱۸۲۱)]

\textbf{ملیت:} ساوویایی (ایتالیایی-فرانسوی)

\textbf{آثار کلیدی:} \book{ملاحظاتی درباره فرانسه}، \book{شب‌های سن‌پترزبورگ}

\textbf{ایده محوری:} ضدانقلاب الهیاتی

\mediumspace

\person{دومستر} رادیکال‌ترین منتقد انقلاب فرانسه بود:

\begin{itemize}[nosep]
    \item انقلاب \textbf{مجازات الهی} برای گناهان فرانسه است
    \item عقل انسانی ناتوان از ساختن نظم است
    \item اقتدار باید \textbf{مطلق و الهی} باشد
    \item پاپ باید رأس نظم اروپایی باشد
    \item حتی \textbf{جلاد} نماد نظم الهی است!
\end{itemize}

\mediumspace

\textbf{نقل‌قول:} \textit{«هر ملتی حکومتی دارد که لایقش است.»}
\end{personbox}

\begin{defbox}[تعریف: ضدروشنگری]
جریانی فکری که:
\begin{itemize}[nosep]
    \item عقل‌گرایی روشنگری را رد می‌کند
    \item پیشرفت را توهم می‌داند
    \item مذهب و سنت را بر علم مقدم می‌شمارد
    \item دموکراسی را آشوب می‌داند
\end{itemize}
نمونه‌ها: دومستر، بونالد، دونوسو کورتس
\end{defbox}

\subsection{رمانتیسم محافظه‌کار}

\begin{historybox}[رمانتیسم سیاسی آلمان]
در آلمان، رمانتیسم با محافظه‌کاری پیوند خورد:
\begin{itemize}[nosep]
    \item نقد عقل‌گرایی و مکانیسم
    \item تأکید بر \textbf{روح ملت} (Volksgeist)
    \item ارزش‌گذاری \textbf{قرون وسطی} و سنت
    \item ضدیت با فردگرایی لیبرال
\end{itemize}

\textbf{چهره‌ها:} نوالیس، آدام مولر، فردریش شلگل

\textbf{میراث:} ناسیونالیسم آلمانی، و متأسفانه برخی عناصر فکری نازیسم
\end{historybox}

% ══════════════════════════════════════════════════════════════════════════════
\section{محافظه‌کاری بریتانیایی}
% ══════════════════════════════════════════════════════════════════════════════

\begin{figure}[htbp]
\centering
\begin{tikzpicture}[scale=0.85]

% عنوان
\node[font=\large\bfseries, text=DarkGray] at (0, 7.5) {\fa{تایم‌لاین: حزب محافظه‌کار بریتانیا}};

% خط زمان
\draw[line width=4pt, MediumGray] (-7, 5.5) -- (7, 5.5);

% رویدادها
% توری‌ها
\fill[RightBlueDark] (-6, 5.5) circle (9pt);
\node[above, text=RightBlueDark, font=\scriptsize\bfseries] at (-6, 6) {\fa{۱۶۸۰s}};
\node[below, text width=2cm, align=center, text=DarkGray, font=\tiny] at (-6, 4.9) {
    \fa{توری‌ها}\\
    \fa{شاه و کلیسا}
};

% پیل
\fill[RightBlue] (-4, 5.5) circle (9pt);
\node[above, text=RightBlue, font=\scriptsize\bfseries] at (-4, 6) {\fa{۱۸۳۴}};
\node[below, text width=2cm, align=center, text=DarkGray, font=\tiny] at (-4, 4.9) {
    \fa{رابرت پیل}\\
    \fa{مانیفست تمورث}
};

% دیزرائیلی
\fill[RightBlueLight] (-1.5, 5.5) circle (10pt);
\node[above, text=RightBlueLight, font=\scriptsize\bfseries] at (-1.5, 6.1) {\fa{۱۸۶۷-۸۰}};
\node[below, text width=2.2cm, align=center, text=DarkGray, font=\tiny] at (-1.5, 4.8) {
    \fa{دیزرائیلی}\\
    \fa{یک‌ملتی}
};

% چرچیل
\fill[RightBlue] (1, 5.5) circle (9pt);
\node[above, text=RightBlue, font=\scriptsize\bfseries] at (1, 6) {\fa{۱۹۴۰-۵۵}};
\node[below, text width=2cm, align=center, text=DarkGray, font=\tiny] at (1, 4.9) {
    \fa{چرچیل}\\
    \fa{جنگ جهانی دوم}
};

% تاچر
\fill[OrangeAccent] (3.5, 5.5) circle (11pt);
\node[above, text=OrangeAccent, font=\scriptsize\bfseries] at (3.5, 6.1) {\fa{۱۹۷۹-۹۰}};
\node[below, text width=2.2cm, align=center, text=DarkGray, font=\tiny] at (3.5, 4.8) {
    \fa{تاچر}\\
    \fa{نئولیبرالیسم}
};

% برگزیت
\fill[GoldAccent] (6, 5.5) circle (10pt);
\node[above, text=GoldAccent, font=\scriptsize\bfseries] at (6, 6) {\fa{۲۰۱۶-۲۰}};
\node[below, text width=2cm, align=center, text=DarkGray, font=\tiny] at (6, 4.8) {
    \fa{برگزیت}\\
    \fa{جانسون}
};

\end{tikzpicture}
\caption{سیر تاریخی محافظه‌کاری بریتانیایی}
\label{fig:tory-timeline}
\end{figure}

\subsection{محافظه‌کاری یک‌ملتی: دیزرائیلی}

\begin{personbox}[بنجامین دیزرائیلی (۱۸۰۴-۱۸۸۱)]

\textbf{ملیت:} بریتانیایی (یهودی‌تبار)

\textbf{سمت:} نخست‌وزیر (۱۸۶۸، ۱۸۷۴-۱۸۸۰)

\textbf{ایده محوری:} محافظه‌کاری یک‌ملتی (One-Nation Conservatism)

\mediumspace

\person{دیزرائیلی} محافظه‌کاری بریتانیایی را بازسازی کرد:

\begin{itemize}[nosep]
    \item بریتانیا به «دو ملت» (ثروتمند و فقیر) تقسیم شده است
    \item اشراف وظیفه دارند از فقرا مراقبت کنند (\textbf{نوبلس اوبلیژ})
    \item اصلاحات اجتماعی برای حفظ نظم لازم است
    \item اتحاد طبقات، نه جنگ طبقاتی
    \item امپراتوری و افتخار ملی
\end{itemize}

\mediumspace

\textbf{دستاوردها:} قانون بهداشت عمومی، مسکن کارگری، گسترش حق رأی

\textbf{نقل‌قول:} \textit{«قدرت به کسانی تعلق دارد که طرف مردم را بگیرند.»}
\end{personbox}

\begin{defbox}[تعریف: محافظه‌کاری یک‌ملتی]
گرایشی در محافظه‌کاری بریتانیایی که:
\begin{itemize}[nosep]
    \item اشرافیت را نه صرفاً امتیاز، بلکه \textbf{مسئولیت} می‌داند
    \item از اصلاحات اجتماعی برای پیشگیری از انقلاب دفاع می‌کند
    \item سرمایه‌داری افسارگسیخته را خطرناک می‌داند
    \item هم ضد سوسیالیسم است، هم ضد لیبرالیسم منچستری
\end{itemize}
\textbf{میراث:} مک‌میلان، هیث، و «محافظه‌کاران مرطوب» (Wets) عصر تاچر
\end{defbox}

\subsection{تاچریسم: انقلاب محافظه‌کارانه؟}

\begin{personbox}[مارگارت تاچر (۱۹۲۵-۲۰۱۳)]

\textbf{ملیت:} بریتانیایی

\textbf{سمت:} نخست‌وزیر (۱۹۷۹-۱۹۹۰)

\textbf{لقب:} «بانوی آهنین»

\mediumspace

\person{تاچر} محافظه‌کاری بریتانیایی را دگرگون کرد—به گونه‌ای که برخی آن را «غیرمحافظه‌کارانه» می‌دانند:

\textbf{سیاست‌ها:}
\begin{itemize}[nosep]
    \item خصوصی‌سازی گسترده
    \item شکستن قدرت اتحادیه‌ها
    \item کاهش مالیات و مقررات‌زدایی
    \item «جامعه وجود ندارد، فقط افراد و خانواده‌ها»
    \item جنگ فالکلند و سخت‌گیری ملی‌گرایانه
\end{itemize}

\mediumspace

\textbf{نقل‌قول معروف:} \textit{«جایگزینی نیست.» (TINA - There Is No Alternative)}

\mediumspace

\textbf{مناقشه:} آیا تاچر محافظه‌کار بود یا لیبرال کلاسیک/نئولیبرال؟
\end{personbox}

\begin{debatebox}[آیا تاچر محافظه‌کار بود؟]
\textbf{بله:}
\begin{itemize}[nosep]
    \item بر خانواده، مذهب، ناسیونالیسم تأکید داشت
    \item ضد سوسیالیسم و اتحادیه‌ها بود
    \item نظم و قانون را مهم می‌دانست
\end{itemize}

\textbf{نه:}
\begin{itemize}[nosep]
    \item محافظه‌کاری سنتی از نهادها دفاع می‌کند—تاچر آن‌ها را ویران کرد
    \item «جامعه وجود ندارد» ضد برک است
    \item رادیکالیسم بازار، محافظه‌کارانه نیست
    \item جوامع محلی و صنایع سنتی را نابود کرد
\end{itemize}

\textbf{شاید:} تاچر ترکیبی بود از محافظه‌کاری اجتماعی + لیبرالیسم اقتصادی.
\end{debatebox}

% ══════════════════════════════════════════════════════════════════════════════
\section{محافظه‌کاری آمریکایی}
% ══════════════════════════════════════════════════════════════════════════════

محافظه‌کاری آمریکایی پدیده‌ای پیچیده است، زیرا آمریکا خود کشوری انقلابی و لیبرال است.

\begin{figure}[htbp]
\centering
\begin{tikzpicture}[scale=0.85]

% عنوان
\node[font=\large\bfseries, text=DarkGray] at (0, 7) {\fa{جناح‌های محافظه‌کاری آمریکایی}};

% دایره‌های همپوشان
% محافظه‌کاری سنتی
\fill[RightBlue, opacity=0.4] (-2.5, 3) circle (2.5cm);
\node[text=RightBlueDark, font=\bfseries] at (-3.5, 4.5) {\fa{سنت‌گرا}};
\node[text=DarkGray, font=\scriptsize, text width=2cm, align=center] at (-3.5, 3) {
    \fa{برک}\\
    \fa{راسل کرک}\\
    \fa{ارزش‌ها}
};

% لیبرتارین
\fill[CenterGreen, opacity=0.4] (2.5, 3) circle (2.5cm);
\node[text=CenterGreenDark, font=\bfseries] at (3.5, 4.5) {\fa{لیبرتارین}};
\node[text=DarkGray, font=\scriptsize, text width=2cm, align=center] at (3.5, 3) {
    \fa{بازار آزاد}\\
    \fa{هایک}\\
    \fa{دولت کوچک}
};

% ضدکمونیست/نئوکان
\fill[OrangeAccent, opacity=0.4] (0, 0.5) circle (2.5cm);
\node[text=OrangeAccent!80!black, font=\bfseries] at (0, -1.2) {\fa{نومحافظه‌کار}};
\node[text=DarkGray, font=\scriptsize, text width=2cm, align=center] at (0, 0.3) {
    \fa{ضدکمونیسم}\\
    \fa{سیاست خارجی}\\
    \fa{قدرت آمریکا}
};

% نقطه اشتراک
\node[text=MediumGray, font=\small\bfseries] at (0, 3) {\fa{فیوژنیسم}};

% راهنما
\node[text=MediumGray, font=\footnotesize, text width=10cm, align=center] at (0, -3) {
    \fa{«فیوژنیسم»: تلاش برای ترکیب این سه جناح در جنبش واحد (از دهه ۱۹۵۰)}
};

\end{tikzpicture}
\caption{سه جناح اصلی محافظه‌کاری آمریکایی}
\label{fig:american-conservatism}
\end{figure}

\subsection{محافظه‌کاری سنتی: راسل کرک}

\begin{personbox}[راسل کرک (۱۹۱۸-۱۹۹۴)]

\textbf{ملیت:} آمریکایی

\textbf{اثر کلیدی:} \book{ذهن محافظه‌کار} (۱۹۵۳)

\textbf{ایده محوری:} احیای سنت برکی در آمریکا

\mediumspace

\person{کرک} تلاش کرد نشان دهد آمریکا هم سنت محافظه‌کاری دارد:

\textbf{ده اصل محافظه‌کاری کرک:}
\begin{enumerate}[nosep]
    \item باور به نظم اخلاقی جاودان
    \item پایبندی به عرف و تداوم
    \item باور به «تجویز» (حق سنتی)
    \item احتیاط در تغییر
    \item احترام به تنوع و تکثر
    \item پذیرش نقصان انسان
    \item پیوند آزادی و مالکیت
    \item دفاع از جوامع داوطلبانه
    \item محدودیت‌های قدرت
    \item ترکیب دگرگونی و حفظ
\end{enumerate}
\end{personbox}

\subsection{ویلیام باکلی و National Review}

\begin{personbox}[ویلیام اف. باکلی جونیور (۱۹۲۵-۲۰۰۸)]

\textbf{ملیت:} آمریکایی

\textbf{نشریه:} \textit{National Review} (تأسیس ۱۹۵۵)

\textbf{ایده محوری:} اتحاد محافظه‌کاران علیه لیبرالیسم

\mediumspace

\person{باکلی} جنبش محافظه‌کار آمریکایی مدرن را ساخت:

\begin{itemize}[nosep]
    \item «ایستادن جلوی تاریخ و فریاد زدن: بایست!»
    \item اتحاد سنت‌گرایان، لیبرتارین‌ها و ضدکمونیست‌ها
    \item طرد افراطیون (جان برچ سوسایتی، نژادپرستان)
    \item آماده‌سازی زمینه برای گلدواتر و ریگان
\end{itemize}
\end{personbox}

\subsection{نومحافظه‌کاری}

\begin{defbox}[تعریف: نومحافظه‌کاری (Neoconservatism)]
جنبشی که از دهه ۱۹۷۰ شکل گرفت:

\textbf{ریشه:} روشنفکران چپ سابق که از چپ ناامید شدند

\textbf{ویژگی‌ها:}
\begin{itemize}[nosep]
    \item سیاست خارجی تهاجمی برای گسترش دموکراسی
    \item ضدیت با شوروی (و بعدها «محور شرارت»)
    \item حمایت از اسرائیل
    \item نقد دولت رفاه (اما نه ضدیت کامل)
    \item نقد فرهنگ دهه ۶۰
\end{itemize}

\textbf{چهره‌ها:} ایروینگ کریستول، نورمن پادورتز، پل ولفوویتز

\textbf{تأثیر:} جنگ عراق ۲۰۰۳
\end{defbox}

\subsection{ریگانیسم}

\begin{personbox}[رونالد ریگان (۱۹۱۱-۲۰۰۴)]

\textbf{ملیت:} آمریکایی

\textbf{سمت:} رئیس‌جمهور (۱۹۸۱-۱۹۸۹)

\textbf{شعار:} «صبح دوباره در آمریکا»

\mediumspace

\person{ریگان} محافظه‌کاری آمریکایی را به قدرت رساند:

\textbf{سیاست‌ها:}
\begin{itemize}[nosep]
    \item کاهش مالیات («ریگانومیکس»)
    \item افزایش هزینه‌های دفاعی
    \item مقررات‌زدایی
    \item ضدکمونیسم قاطع
    \item اتحاد با راست مذهبی
\end{itemize}

\mediumspace

\textbf{نقل‌قول:} \textit{«دولت راه‌حل مشکل ما نیست؛ دولت خود مشکل است.»}
\end{personbox}

\begin{figure}[htbp]
\centering
\begin{tikzpicture}[scale=0.85]

% عنوان
\node[font=\large\bfseries, text=DarkGray] at (0, 6.5) {\fa{ائتلاف ریگان: سه پایه}};

% سه پایه
% اقتصادی
\fill[RightBlue, opacity=0.7, rounded corners=8pt] (-5, 1) rectangle (-1.5, 5);
\draw[line width=2pt, color=RightBlue, rounded corners=8pt] (-5, 1) rectangle (-1.5, 5);
\node[text=white, font=\bfseries] at (-3.25, 4.5) {\fa{اقتصادی}};
\node[text=VeryLightGray, font=\scriptsize, text width=3cm, align=center] at (-3.25, 2.8) {
    \fa{کاهش مالیات}\\
    \fa{مقررات‌زدایی}\\
    \fa{بازار آزاد}\\
    \fa{لیبرتارین‌ها}
};

% اجتماعی/مذهبی
\fill[PurpleAccent, opacity=0.7, rounded corners=8pt] (-1, 1) rectangle (2.5, 5);
\draw[line width=2pt, color=PurpleAccent, rounded corners=8pt] (-1, 1) rectangle (2.5, 5);
\node[text=white, font=\bfseries] at (0.75, 4.5) {\fa{اجتماعی}};
\node[text=VeryLightGray, font=\scriptsize, text width=3cm, align=center] at (0.75, 2.8) {
    \fa{ضد سقط جنین}\\
    \fa{ارزش‌های خانواده}\\
    \fa{دعا در مدرسه}\\
    \fa{راست مذهبی}
};

% سیاست خارجی
\fill[OrangeAccent, opacity=0.7, rounded corners=8pt] (3, 1) rectangle (6.5, 5);
\draw[line width=2pt, color=OrangeAccent, rounded corners=8pt] (3, 1) rectangle (6.5, 5);
\node[text=white, font=\bfseries] at (4.75, 4.5) {\fa{امنیتی}};
\node[text=VeryLightGray, font=\scriptsize, text width=3cm, align=center] at (4.75, 2.8) {
    \fa{ضدکمونیسم}\\
    \fa{قدرت نظامی}\\
    \fa{ناسیونالیسم}\\
    \fa{نومحافظه‌کاران}
};

% نقطه اتصال
\fill[GoldAccent] (0.75, -0.3) circle (10pt);
\node[text=white, font=\bfseries\small] at (0.75, -0.3) {\fa{ریگان}};

% خطوط اتصال
\draw[line width=2pt, color=GoldAccent] (-3.25, 1) -- (0.75, -0.1);
\draw[line width=2pt, color=GoldAccent] (0.75, 1) -- (0.75, -0.1);
\draw[line width=2pt, color=GoldAccent] (4.75, 1) -- (0.75, -0.1);

\end{tikzpicture}
\caption{ائتلاف سه‌گانه ریگان}
\label{fig:reagan-coalition}
\end{figure}

% ══════════════════════════════════════════════════════════════════════════════
\section{محافظه‌کاری اجتماعی}
% ══════════════════════════════════════════════════════════════════════════════

\begin{defbox}[تعریف: محافظه‌کاری اجتماعی]
گرایشی که بر حفظ ارزش‌ها و نهادهای اجتماعی سنتی تأکید دارد:
\begin{itemize}[nosep]
    \item \textbf{خانواده:} خانواده هسته‌ای، نقش‌های سنتی جنسیتی
    \item \textbf{مذهب:} اخلاق مبتنی بر دین، حضور دین در عرصه عمومی
    \item \textbf{جامعه محلی:} پیوندهای محلی و داوطلبانه
    \item \textbf{فرهنگ:} دفاع از فرهنگ ملی/غربی
    \item \textbf{نظم:} قانون و نظم، تنبیه جرم
\end{itemize}
\end{defbox}

\begin{historybox}[راست مذهبی در آمریکا]
از دهه ۱۹۷۰، مسیحیان انجیلی وارد سیاست شدند:

\textbf{مسائل کلیدی:}
\begin{itemize}[nosep]
    \item سقط جنین (پس از حکم رو در برابر وید، ۱۹۷۳)
    \item دعا در مدارس دولتی
    \item ازدواج همجنس‌گرایان
    \item آموزش جنسی
    \item تکامل‌گرایی در مدارس
\end{itemize}

\textbf{سازمان‌ها:} اکثریت اخلاقی (جری فالول)، ائتلاف مسیحی (پت رابرتسون)

\textbf{تأثیر:} پایگاه انتخاباتی کلیدی جمهوری‌خواهان
\end{historybox}

% ══════════════════════════════════════════════════════════════════════════════
\section{محافظه‌کاری و دولت رفاه}
% ══════════════════════════════════════════════════════════════════════════════

رابطه محافظه‌کاری با دولت رفاه پیچیده است. برخی محافظه‌کاران از آن دفاع کرده‌اند، برخی مخالفند.

\subsection{محافظه‌کاری پدرسالارانه: بیسمارک}

\begin{personbox}[اتو فون بیسمارک (۱۸۱۵-۱۸۹۸)]

\textbf{ملیت:} آلمانی (پروسی)

\textbf{سمت:} صدراعظم آلمان (۱۸۷۱-۱۸۹۰)

\textbf{لقب:} «صدراعظم آهنین»

\mediumspace

\person{بیسمارک} بنیان‌گذار دولت رفاه مدرن بود—نه از روی سوسیالیسم، بلکه برای مقابله با آن:

\begin{itemize}[nosep]
    \item بیمه بیماری (۱۸۸۳)
    \item بیمه حوادث کار (۱۸۸۴)
    \item بیمه بازنشستگی (۱۸۸۹)
\end{itemize}

\textbf{منطق:} با دادن امنیت به کارگران، آن‌ها را از سوسیالیست‌ها دور کنیم.

\textbf{نقل‌قول:} \textit{«سوسیالیسم دولتی من را قبول کنید، وگرنه سوسیالیسم انقلابی خواهید داشت.»}
\end{personbox}

\begin{comparebox}[مقایسه: محافظه‌کاران و دولت رفاه]
\begin{table}[H]
\centering
\small
\begin{tabular}{r|c|c}
\toprule
\textbf{گرایش} & \textbf{موضع} & \textbf{نماینده} \\
\midrule
پدرسالارانه & دولت رفاه لازم است & بیسمارک، دیزرائیلی \\
یک‌ملتی & برای انسجام اجتماعی ضروری & مک‌میلان، مرکل \\
لیبرتارین & مخالف؛ تجاوز به آزادی & هایک، راتبارد \\
نئولیبرال & کاهش و خصوصی‌سازی & تاچر، ریگان \\
ملی‌گرای جدید & حفظ برای «خودی‌ها» & اوربان، ترامپ \\
\bottomrule
\end{tabular}
\end{table}
\end{comparebox}

\subsection{Red Tory: محافظه‌کاری اجتماعی جدید}

\begin{defbox}[تعریف: Red Tory]
گرایشی در محافظه‌کاری بریتانیایی و کانادایی که می‌کوشد:
\begin{itemize}[nosep]
    \item هم با نئولیبرالیسم (بازار افسارگسیخته) مخالفت کند
    \item هم با دولت بوروکراتیک
    \item تأکید بر جوامع محلی، تعاونی‌ها، نهادهای مدنی
    \item «توزیع‌گرایی» (Distributism): مالکیت گسترده به جای تمرکز
\end{itemize}
\textbf{چهره:} فیلیپ بلاند (\book{Red Tory}، ۲۰۱۰)
\end{defbox}

% ══════════════════════════════════════════════════════════════════════════════
\section{اختلافات درونی محافظه‌کاری}
% ══════════════════════════════════════════════════════════════════════════════

\begin{figure}[htbp]
\centering
\begin{tikzpicture}[scale=0.85]

% عنوان
\node[font=\large\bfseries, text=DarkGray] at (0, 6.5) {\fa{محورهای اختلاف درون محافظه‌کاری}};

% محور افقی: سنت vs بازار
\draw[line width=3pt, MediumGray, -{Stealth[length=8pt]}] (-6, 3) -- (6, 3);
\node[text=PurpleAccent, font=\small\bfseries] at (-6.5, 3) {\fa{سنت}};
\node[text=CenterGreen, font=\small\bfseries] at (6.5, 3) {\fa{بازار}};

% محور عمودی: ملی‌گرا vs جهان‌وطن
\draw[line width=3pt, MediumGray, -{Stealth[length=8pt]}] (0, 0) -- (0, 6);
\node[text=OrangeAccent, font=\small\bfseries, rotate=90] at (-0.5, 0) {\fa{جهان‌وطن}};
\node[text=LeftRed, font=\small\bfseries, rotate=90] at (-0.5, 5.8) {\fa{ملی‌گرا}};

% نقاط
% محافظه‌کار سنتی
\fill[PurpleAccent] (-4, 4.5) circle (8pt);
\node[text=DarkGray, font=\scriptsize] at (-4, 3.9) {\fa{سنت‌گرا}};

% لیبرتارین
\fill[CenterGreen] (4, 1.5) circle (8pt);
\node[text=DarkGray, font=\scriptsize] at (4, 0.9) {\fa{لیبرتارین}};

% نئوکان
\fill[OrangeAccent] (2, 1.5) circle (8pt);
\node[text=DarkGray, font=\scriptsize] at (2, 0.9) {\fa{نئوکان}};

% ملی‌گرای جدید
\fill[LeftRed] (-2, 5) circle (8pt);
\node[text=DarkGray, font=\scriptsize] at (-2, 5.6) {\fa{ملی‌گرای جدید}};

% نئولیبرال
\fill[RightBlue] (3, 3.5) circle (8pt);
\node[text=DarkGray, font=\scriptsize] at (3, 4.1) {\fa{نئولیبرال}};

\end{tikzpicture}
\caption{محورهای اختلاف درون محافظه‌کاری}
\label{fig:conservatism-divisions}
\end{figure}

\begin{debatebox}[اختلاف کلیدی: سنت یا بازار؟]
\textbf{محافظه‌کاران سنتی:}
\begin{itemize}[nosep]
    \item سرمایه‌داری افسارگسیخته جوامع را ویران می‌کند
    \item بازار باید در خدمت جامعه باشد، نه برعکس
    \item برخی چیزها نباید کالایی شوند (بهداشت، آموزش؟)
    \item نقد «خلاقیت ویرانگر» سرمایه‌داری
\end{itemize}

\textbf{محافظه‌کاران بازارمحور (لیبرتارین/نئولیبرال):}
\begin{itemize}[nosep]
    \item بازار آزاد، بهترین تضمین آزادی است
    \item دولت بزرگ، تهدید اصلی است
    \item سنت‌ها باید با رقابت بازار محک بخورند
    \item «محافظه‌کاری» یعنی حفظ آزادی، نه نهادهای کهنه
\end{itemize}

\textbf{این اختلاف در عصر ترامپ و برگزیت به اوج رسید.}
\end{debatebox}

\begin{debatebox}[اختلاف دیگر: ملی‌گرایی یا جهان‌وطنی؟]
\textbf{ملی‌گرایان:}
\begin{itemize}[nosep]
    \item ملت واحد اصلی سیاست است
    \item مرزها و هویت ملی مهم‌اند
    \item جهانی‌شدن به ضرر کارگران بومی است
    \item «ابتدا آمریکا» / «ابتدا بریتانیا»
\end{itemize}

\textbf{جهان‌وطن‌ها:}
\begin{itemize}[nosep]
    \item تجارت آزاد به نفع همه است
    \item ناسیونالیسم افراطی خطرناک است
    \item غرب باید متحد بماند
    \item مهاجرت برای اقتصاد لازم است
\end{itemize}
\end{debatebox}

% ══════════════════════════════════════════════════════════════════════════════
\section{نقد محافظه‌کاری}
% ══════════════════════════════════════════════════════════════════════════════

\subsection{نقد از چپ}

\begin{itemize}[nosep, rightmargin=1.5cm]
    \item محافظه‌کاری \textbf{توجیه نابرابری} است—دفاع از منافع ثروتمندان
    \item «سنت» یعنی ادامه ستم‌های تاریخی (برده‌داری، پدرسالاری)
    \item «احتیاط» یعنی جلوگیری از پیشرفت
    \item محافظه‌کاران همیشه در طرف اشتباه تاریخ بوده‌اند
\end{itemize}

\subsection{نقد از لیبرال‌ها}

\begin{itemize}[nosep, rightmargin=1.5cm]
    \item محافظه‌کاری اجتماعی، آزادی فردی را محدود می‌کند
    \item «ارزش‌های سنتی» یعنی کنترل زنان و اقلیت‌ها
    \item چرا باید به سنت احترام گذاشت اگر ناعادلانه است؟
\end{itemize}

\subsection{خودانتقادی: چه چیزی ارزش حفظ دارد؟}

\begin{quotebox}[راجر اسکروتن، فیلسوف محافظه‌کار]
\textit{«محافظه‌کار واقعی می‌داند که برخی چیزها ارزش حفظ کردن ندارند. هنر محافظه‌کاری، تشخیص این است که چه چیزی واقعاً ارزشمند است.»}
\end{quotebox}

% ══════════════════════════════════════════════════════════════════════════════
\section{محافظه‌کاری معاصر}
% ══════════════════════════════════════════════════════════════════════════════

\begin{historybox}[بحران محافظه‌کاری معاصر]
پس از تاچر و ریگان، محافظه‌کاری دچار بحران هویت شده:
\begin{itemize}[nosep]
    \item جنگ عراق نومحافظه‌کاری را بی‌اعتبار کرد
    \item بحران مالی ۲۰۰۸ نئولیبرالیسم را زیر سؤال برد
    \item جهانی‌شدن طبقه کارگر سفید را به سمت پوپولیسم برد
    \item پوپولیسم ملی‌گرا (ترامپ، برگزیت) نظم سنتی را برهم زد
\end{itemize}

\textbf{سؤال:} محافظه‌کاری پس از ترامپ چه خواهد بود؟
\end{debatebox}

\begin{table}[htbp]
\centering
\caption{انواع محافظه‌کاری معاصر}
\label{tab:contemporary-conservatism}
\begin{tabular}{r|r|r}
\toprule
\textbf{نوع} & \textbf{ویژگی کلیدی} & \textbf{نماینده} \\
\midrule
سنت‌گرای کلاسیک & برک، کرک، ارزش‌ها & راجر اسکروتن \\
نئولیبرال & بازار آزاد، دولت کوچک & (تاچر، ریگان) \\
نومحافظه‌کار & سیاست خارجی تهاجمی & (چنی، ولفوویتز) \\
ملی‌گرای پوپولیست & ضد نخبگان، ضد مهاجرت & ترامپ، اوربان \\
محافظه‌کار ملی جدید & دولت فعال برای ملت & (مارکو روبیو؟) \\
\bottomrule
\end{tabular}
\end{table}

% ══════════════════════════════════════════════════════════════════════════════
%                         خلاصه و پرسش‌ها
% ══════════════════════════════════════════════════════════════════════════════

\newpage

\begin{summarybox}

\textbf{۱. تعریف:} محافظه‌کاری هم به معنای مخالفت با تغییر رادیکال است، هم دفاع از ارزش‌های خاص (سنت، نظم، خانواده، مذهب).

\textbf{۲. برک:} پدر محافظه‌کاری مدرن. خرد جمعی سنت، جامعه ارگانیک، اصلاح تدریجی.

\textbf{۳. قاره‌ای:} ضدانقلابیون (دومستر)، رمانتیسم محافظه‌کار آلمانی.

\textbf{۴. بریتانیا:} از توری‌ها تا تاچر. دیزرائیلی (یک‌ملتی) vs تاچر (نئولیبرال).

\textbf{۵. آمریکا:} سه جناح: سنت‌گرا، لیبرتارین، نومحافظه‌کار. فیوژنیسم و ائتلاف ریگان.

\textbf{۶. اجتماعی:} خانواده، مذهب، ارزش‌های سنتی. راست مذهبی.

\textbf{۷. دولت رفاه:} برخی محافظه‌کاران (بیسمارک، یک‌ملتی) موافق؛ لیبرتارین‌ها مخالف.

\textbf{۸. اختلافات:} سنت vs بازار؛ ملی‌گرایی vs جهان‌وطنی.

\textbf{۹. امروز:} بحران هویت پس از نئولیبرالیسم. پوپولیسم ملی‌گرا چالش جدید.

\end{summarybox}

\vspace{20pt}

\begin{questionbox}

\begin{enumerate}[nosep, rightmargin=1cm]
    \item اگر محافظه‌کاری یعنی حفظ وضع موجود، آیا محافظه‌کار سوسیالیست هم ممکن است؟ (کسی که می‌خواهد دولت رفاه را «حفظ» کند)

    \item آیا تاچر و ریگان واقعاً «محافظه‌کار» بودند؟ یا رادیکال‌هایی بودند که جامعه را دگرگون کردند؟

    \item برک از انقلاب آمریکا دفاع کرد اما انقلاب فرانسه را محکوم کرد. آیا این تمایز موجه است؟

    \item آیا می‌توان هم محافظه‌کار اجتماعی بود (ارزش‌های سنتی) و هم لیبرال اقتصادی (بازار آزاد)؟ یا این دو در تضادند؟
\end{enumerate}

\end{questionbox}

% ══════════════════════════════════════════════════════════════════════════════
%                         پایان فصل نهم
% ══════════════════════════════════════════════════════════════════════════════